\section{Introduction}
An upstream renewal desk can know an accepted obligation that a downstream execution gateway does not retain. When only a small service symbol crosses that boundary, the provider must choose both its terminal instruction alphabet and the branch-dependent encoder. This is a quantizer-design problem, but its objective and feasible randomization are generated by a contract: an aggregate continuation promise, local participation limits, and the cost of intermediate rather than terminal service. We ask how to choose the alphabet and allocate the promise jointly when individual lottery realizations have an overrun limit and enabling an additional execution level has a cost.

Our algorithmic results use this contractual structure rather than enumerating all its possible local decisions. After assigning a price to the aggregate promise, an optimal ordered alphabet can be represented by a path whose edges account for the branch caps they cross. A dynamic program solves this price problem in polynomial time, with both the number of branches and the alphabet budget allowed to vary. The same decomposition gives exact bounds for partially chosen alphabets. A finite global search then returns either an exact catalog optimum or an explicit lower and upper bound at interruption. Its witness records the complete search partition and can be checked independently of the search code. No equality between the price relaxation and the original problem is presumed.

The new recovery result makes this price oracle constructive even when the original price gap is positive. Two priced solutions determine a single installed union alphabet and exact branch targets; canonical union lotteries preserve the pre-draw service rule and every realization ceiling. With zero opening charges, allowing at most twice as many execution levels gives any prescribed additive accuracy relative to the original-budget catalog optimum, in polynomial computation with variable branch count and alphabet budget. Nonnegative opening charges enter through an explicit excess-charge term. This is a quantified resource tradeoff, not randomization of the installed interface or a claim that extra symbols are free. An exact example exhibits a strictly positive original price gap and its removal by the fixed union construction.

The method is part of a unified theory covering every feasible aggregate promise, continuously chosen execution levels, bounded realization exposure, and charges on the actual selected levels. A branch-specific eligibility prefix and a canonical adjacent lottery reduce each fixed-book problem to concave resource allocation. The general concave model admits this exact fixed-book structure, the priced decomposition, and a feasible discretization bound; exact unrestricted continuous optimization is proved for the rational quadratic subclass. For that subclass, eliminating lottery probabilities gives closed quadratic cells, and exhaustive stationary-face search supplies a finite exact reference algorithm, including collisions and nonsaturated pathwise participation. We retain that result and state its full exponential cost.

A second connection turns a numerical catalog into a controlled approximation of continuous design. Downward rounding and reconstructed lottery probabilities preserve every branch promise and realization ceiling exactly, even at binding constraints. A uniform error bound converts the search interval into a continuous-value interval, separating incomplete optimization from discretization. We prove that its linear mesh rate is order-sharp on a single fixed, uncharged input. The bound applies before catalog search finishes. Combined with the new recovery theorem, it also gives a polynomial continuous-design comparison with an expanded alphabet; the unexpanded fixed-budget enumeration guarantee remains a separate result.

The saturated contract remains especially tractable. Deterministic execution gives an ordered one-sided quantizer. Under expected participation, a quadratic optimum has cap-anchored nonhighest levels and a continuously optimized highest level that can lie between caps. A crossing-difference identity preserves Monge order after continuous tail minimization, allowing classical linear work per dynamic-programming layer. We retain those results and proofs, the all-promise exact-alphabet phase diagram, stability bounds, and the faster saturated charged-catalog recurrence. The new decomposition broadens the algorithmic treatment without replacing these sharper endpoint results.

Ordered quantizer dynamic programming, optimized reproduction points, unbiased splitting, Lagrangian relaxation, and matrix search are established ideas \citep{Wu1991,PagesWilbertz2012,CrociEtAl2022,JourdainPages2021,HandlerZang1980,Lemarechal2001}. The distinct result is how exact promises, heterogeneous shortfall costs, support eligibility, and a shared selected alphabet produce a polynomial priced path problem and reusable completion bounds. We formulate this result for an ordered-interface allocation class that does not require the renewal interpretation. The operational model remains an explicitly stipulated risk-neutral or actuarial acceptance rule followed by restricted execution; it is not presented as customer data or an implemented field deployment. Synthetic studies test exact bounds, expose the continuous reference solver's enumeration frontier, and compare a separate original-policy nonlinear formulation. A monotone-charge theorem also identifies precisely why the retained one-branch randomization example depends on its specially shaped nonmonotone charge.
